%%% math-article.tex
%%% Template for pure mathematics articles using AMS-LaTeX
%%% Based on AMS best practices and community conventions
%%%------------------------------------------------------------------
\documentclass[11pt]{amsart}

%%% Packages
\usepackage{mathtools}   % loads amsmath, fixes bugs, adds extensions
\usepackage{amssymb}     % \mathbb, \mathfrak, extra symbols
\usepackage{bm}          % bold math symbols (\bm)
\usepackage{graphicx}    % includegraphics (uncomment if needed)
% \usepackage{hyperref}  % hyperlinks (uncomment if needed, load last)

%%% Equation numbering
\numberwithin{equation}{section}

%%% Theorem environments
% plain style: italic body (theorems, lemmas, propositions, corollaries)
\theoremstyle{plain}
\newtheorem{thm}{Theorem}[section]
\newtheorem{lem}[thm]{Lemma}
\newtheorem{prop}[thm]{Proposition}
\newtheorem{cor}[thm]{Corollary}

% definition style: upright body (definitions, examples)
\theoremstyle{definition}
\newtheorem{defn}[thm]{Definition}
\newtheorem{ex}[thm]{Example}

% remark style: upright body (remarks, notes)
\theoremstyle{remark}
\newtheorem{rem}[thm]{Remark}
\newtheorem*{note}{Note}

%%% Math operators (define your own)
\DeclareMathOperator{\sgn}{sgn}
\DeclareMathOperator{\id}{id}
\DeclareMathOperator{\im}{im}
\DeclareMathOperator{\Hom}{Hom}
\DeclareMathOperator{\Ext}{Ext}
\DeclareMathOperator{\Tor}{Tor}
% \DeclareMathOperator*{\limsup}{lim\,sup}  % limits below (like \lim)

%%% Common macros — number sets
\newcommand{\R}{\mathbb{R}}
\newcommand{\C}{\mathbb{C}}
\newcommand{\Z}{\mathbb{Z}}
\newcommand{\N}{\mathbb{N}}
\newcommand{\Q}{\mathbb{Q}}
\newcommand{\F}{\mathbb{F}}

%%% Common macros — bold vectors/matrices
\newcommand{\vect}[1]{\bm{#1}}
\newcommand{\mat}[1]{\bm{#1}}

%%% Common macros — differential (choose ONE convention)
% Convention 1: upright d (ISO standard, common in applied math)
\newcommand{\dd}{\mathrm{d}}
% Convention 2: italic d with thin space (pure math convention, Knuth)
% \newcommand{\dd}{\,d}

%%% Common macros — formatting
\newcommand{\norm}[1]{\lVert #1 \rVert}
\newcommand{\abs}[1]{\lvert #1 \rvert}
\newcommand{\inner}[2]{\langle #1, #2 \rangle}
\newcommand{\set}[1]{\{\, #1 \,\}}

%%%------------------------------------------------------------------
\begin{document}

\title[Short title]{Full Title of the Paper}
\author{Author Name}
\address{Department of Mathematics, University, City, Country}
\email{author@university.edu}
\date{\today}

\subjclass[2020]{Primary 00A00; Secondary 00B00}
\keywords{keyword one, keyword two, keyword three}

\begin{abstract}
Write the abstract here. Summarize the main results and techniques.
\end{abstract}

\maketitle

%%%------------------------------------------------------------------
\section{Introduction}
\label{sec:intro}

Write the introduction here. Use inline math like $f \colon X \to Y$.

A single numbered equation:
\begin{equation}
  \label{eq:example}
  \int_0^1 f(x) \, \dd x = F(1) - F(0).
\end{equation}

Reference it with \eqref{eq:example}.

%%%------------------------------------------------------------------
\section{Main Results}
\label{sec:main}

\begin{thm}[Main Theorem]
  \label{thm:main}
  Let $X$ be a complete metric space and $T \colon X \to X$ a contraction
  mapping with constant $k < 1$. Then $T$ has a unique fixed point $x^* \in X$.
\end{thm}

\begin{proof}
  Define the sequence $x_{n+1} = T(x_n)$ for some starting point $x_0 \in X$.
  We show $\{x_n\}$ is Cauchy:
  \begin{align*}
    \norm{x_{n+1} - x_n}
      &= \norm{T(x_n) - T(x_{n-1})} \\
      &\le k \norm{x_n - x_{n-1}} \\
      &\le k^n \norm{x_1 - x_0}.
  \end{align*}
  Since $k < 1$, the geometric series converges, so $\{x_n\}$ is Cauchy.
  Completeness gives convergence; continuity of $T$ gives the fixed point.
\end{proof}

\begin{cor}
  \label{cor:apply}
  The fixed point depends continuously on the contraction map $T$.
\end{cor}

\begin{defn}
  \label{defn:lipschitz}
  A function $f \colon X \to Y$ between metric spaces is \emph{Lipschitz}
  if there exists $L \ge 0$ such that $\norm{f(x) - f(y)} \le L \norm{x - y}$
  for all $x, y \in X$.
\end{defn}

\begin{rem}
  Every contraction is Lipschitz with constant $L = k < 1$.
\end{rem}

%%% Aligned equations with shared number
\begin{equation}
  \label{eq:multiline}
  \begin{split}
    a &= b + c \\
      &= d + e + f \\
      &= g.
  \end{split}
\end{equation}

%%% Gathered equations (centered, each numbered)
\begin{gather}
  \label{eq:gather1}
  \alpha = \beta + \gamma, \\
  \label{eq:gather2}
  \delta = \epsilon.
\end{gather}

%%%------------------------------------------------------------------
\section{Conclusion}
\label{sec:conclusion}

Summarize results and discuss future directions.

%%%------------------------------------------------------------------
%%% Bibliography (bibtex + amsplain)
%%% Create a .bib file and uncomment:
% \bibliographystyle{amsplain}
% \bibliography{references}

%%% Or inline bibliography (amsrefs style):
% \begin{bibdiv}
%   \begin{biblist}
%     \bib{key}{article}{
%       author={Name, A.},
%       title={Title},
%       journal={Journal},
%       date={2024},
%       volume={1},
%       pages={1--10},
%     }
%   \end{biblist}
% \end{bibdiv}

\end{document}
